%========================================================
% D2 — Partial Trace for Maximally Entangled States
%========================================================

\section{Partial Trace for Maximally Entangled States}

This derivation establishes explicitly that the reduced density operators of a maximally entangled Bell state are maximally mixed.

\medskip

Consider the Bell state

\begin{equation}
|\Psi^+\rangle
=
\frac{1}{\sqrt{2}}
\left(
|01\rangle + |10\rangle
\right).
\label{eq:bell_state_psiplus_derivation}
\end{equation}

The corresponding density operator is

\begin{equation}
\rho_{AB}
=
|\Psi^+\rangle \langle \Psi^+|.
\label{eq:bell_density_operator_derivation}
\end{equation}

Expanding explicitly:

\begin{align}
\rho_{AB}
&=
\frac{1}{2}
\left(
|01\rangle + |10\rangle
\right)
\left(
\langle 01| + \langle 10|
\right)
\nonumber
\\
&=
\frac{1}{2}
\Big(
|01\rangle \langle 01|
+
|01\rangle \langle 10|
+
|10\rangle \langle 01|
+
|10\rangle \langle 10|
\Big).
\label{eq:bell_density_expansion}
\end{align}

\medskip

The reduced density operator of subsystem \( A \) is obtained by tracing over subsystem \( B \):

\begin{equation}
\rho_A
=
\mathrm{Tr}_B(\rho_{AB}).
\label{eq:partial_trace_definition_derivation}
\end{equation}

Using the computational basis \( \{|0\rangle, |1\rangle\} \) for subsystem \( B \), the partial trace can be written as

\begin{equation}
\rho_A
=
\sum_{b=0}^{1}
(I_A \otimes \langle b|)
\rho_{AB}
(I_A \otimes |b\rangle).
\label{eq:partial_trace_computational_basis}
\end{equation}

Each contribution in Eq.~\eqref{eq:bell_density_expansion} can now be evaluated separately.

\medskip

For the first diagonal term:

\begin{align}
\mathrm{Tr}_B
\left(
|01\rangle \langle 01|
\right)
&=
|0\rangle \langle 0|
\,
\langle 1|1\rangle
\nonumber
\\
&=
|0\rangle \langle 0|.
\end{align}

Similarly,

\begin{equation}
\mathrm{Tr}_B
\left(
|10\rangle \langle 10|
\right)
=
|1\rangle \langle 1|.
\end{equation}

The off-diagonal terms vanish because of basis orthogonality:

\begin{align}
\mathrm{Tr}_B
\left(
|01\rangle \langle 10|
\right)
&=
|0\rangle \langle 1|
\,
\langle 1|0\rangle
\nonumber
\\
&=
0,
\end{align}

and analogously,

\begin{equation}
\mathrm{Tr}_B
\left(
|10\rangle \langle 01|
\right)
=
0.
\end{equation}

Collecting all contributions gives

\begin{align}
\rho_A
&=
\frac{1}{2}
\left(
|0\rangle \langle 0|
+
|1\rangle \langle 1|
\right)
\nonumber
\\
&=
\frac{1}{2} I.
\label{eq:reduced_density_maximally_mixed}
\end{align}

By symmetry, the same result holds for subsystem \( B \):

\begin{equation}
\rho_B
=
\frac{1}{2} I.
\label{eq:reduced_density_B_maximally_mixed}
\end{equation}

\medskip

This result has a direct physical interpretation.

Although the global bipartite state \( |\Psi^+\rangle \) is pure, each subsystem individually is described by a maximally mixed state.

Consequently, no local measurement performed on a single subsystem can reveal complete information about the global quantum state.

The information characterizing the entangled state is therefore encoded entirely in the nonlocal correlations between the two subsystems.

\medskip

This property constitutes one of the fundamental signatures of quantum entanglement and establishes the conceptual connection between entanglement and local mixedness explored throughout the present work.
```
